
\subsubsection{4.1 Research Questions}

\textbf{RQ1:} Does epoch-5 k-means clustering recover spurious group structure on a dataset where the spurious attribute is a dominant, scene-level feature aligned with the ImageNet pretraining prior (Waterbirds)?

\textbf{RQ2:} Does the same pipeline succeed on a dataset where the spurious attribute is a fine-grained texture feature not well-represented in the pretraining prior (CelebA)?

\textbf{RQ3:} Is any observed performance gap attributable to feature-type differences rather than to experimental confounds?

\subsubsection{4.2 Baselines}

A random clustering baseline is computed by randomly shuffling cluster assignments while preserving the cluster-size distribution, then computing AMI and worst-cluster purity against ground-truth group labels. Results: Waterbirds random AMI ≈ 0.0001, random purity = 0.728; CelebA random AMI ≈ 0.000, random purity = 0.440.

Note that the random purity for Waterbirds (0.728) is above the threshold of 0.75 for purity alone, which reflects the fact that random assignment into k=2 clusters on a 4-group dataset with unequal group sizes produces a non-trivial baseline purity. The AMI threshold (≥ 0.5) is the more informative criterion under this setup.

\subsubsection{4.3 Evaluation Metrics}

\textbf{Adjusted Mutual Information (AMI)} measures information-theoretic agreement between cluster assignments and ground-truth group labels, adjusted for chance. It is computed using sklearn.metrics.adjusted_mutual_info_score. AMI=0 for random performance; AMI=1 for perfect alignment. The pre-specified success threshold is AMI ≥ 0.5.

\textbf{Worst-Cluster Purity} measures the minimum fraction of the dominant group within any single cluster, i.e., min_c (max_g count(g, c) / count(c)) over clusters c and groups g. The pre-specified success threshold is purity ≥ 0.75.

Both metrics are computed against ground-truth group labels that are not used during training or clustering.

